\section{The Idea of the World in the East}

This very early review by \textbf{Julius Evola} of ``The Science of Peace'' and ``The Science of the Sacred Word'' originally appeared in \textit{Bilychnis}, volume XXVI, September, 1925. This is part one of two parts. The third part will deal with the idea of the world in the West.\footnote{This is not the American musician who goes by that name.}

\begin{quotationx}
As readers here know, Evola's first book, The Individual and the Becoming of the World, was presented to the Free Theosophical Society in Rome; hence, he was aware of theosophical literature. \textbf{Bhagavan Das} was one of the figures who stood out to Evola, although apparently ignored by the Theosophical Society today. Evola expresses the reasons below, but it should be clear that Bhagavan Das' knowledge of German idealism, particularly the thought of Fichte made an impression on Evola.

Bhagavan Das was an Indian Theosophist writing in the first half of the twentieth century. In true theosophical fashion, he claimed to have uncovered the Pranava Vada by \textbf{Gargayana}, which was dictated to him by a blind man who had memorized the million lines. Nevertheless, Evola has extraordinary high praise for the man and his works, claiming it was written for the Western mind. That is no doubt a side effect of his theosophical training, but he also relates Vedic ideas to Western thought, especially the German idealists.

In my research into Bhagavan Das, I found out that his works were also highly recommended by the Theosophist \textbf{Alice Bailey}. Curiously, there are some pagan circles in Russia\footnote{\url{https://web.archive.org/web/20170711193147/http://sicsa.huji.ac.il:80/13shnir.html}} that hold both Evola and Bailey both in high esteem, for reasons that escape me.

For those interested in the texts:

\textit{The Science of Peace:} \url{https://gornahoor.net/library/ScienceOfPeace.pdf}

\textit{The Science of the Sacred Word:} \url{http://www.makara.us/05ref/01books/pranavavada/pv\_toc.htm} 

\textit{Background information on the Pranava Vada:} \url{http://thecelestialnavigator.blogspot.com/2010/12/pranava-time.html}
\end{quotationx}

\subsection*{First part}

There is a point that whoever approaches oriental wisdom should keep constantly in mind and that instead, as far as I know, is neglected also by our major philologists is this: the quality and role by which its expressions in the East and in the West is experienced are in fact different. While among us, to experience a content means, approximately, to think it and to know how to speak it, so that the expression in its rational form is something that expresses it adequately and with which it is interchangeable, in the Oriental, the content instead draws its evidence from an interior and individual experience and, in respect to which, its expression is always an allusive sign, a symbol, a contingent image. Hence, simply translated oriental texts imported into Europe can only yield little or nothing: to translate an oriental text should not mean to put our words in the place, e.g., of Sanskrit words, but rather to seek for an inner reconstruction to take into account the profound explanation of that content, which in the Orient is supported not on the expression, but rather on an internal apperception, and therefore to make oneself give this same content back to rational mediation according to its own form. Such is the true translation, in respect to which the linguistic translation is something in fact secondary and subordinate, and only in this conditions can oriental metaphysics fit into the western mentality and can be not a simple object of erudite curiosity, but can operate in a vital way.

Now the work of Bhagavan Das, which we highly recommend to the reader, represents the most serious and successful effort along such lines. Bhagavan Das is an Indian who has sketched out an austere discipline of western speculative culture, and joining it with singular gifts of penetration and critical awareness, he has proposed to give the essence of one of the greatest Indian metaphysical systems---the Vedanta---in an explanation proper to our European mentality, and thus (and in that is the superiority of his work in respect to various purely intellectual presentations previously attempted) to the dialectical and critical awareness elaborated by German idealism, in which that mentality has culminated. The Science of Peace represents a general placement, a taking of a purely speculative, synthetic and fully detailed position; the redaction of Gargayana's Pranava Vada is its later development that is deployed and shares in the opulent material of the Vedas, thereby leading the reader into the heart of the Indian tradition. It treats a work so complex and rich that the attempt to give it here in few words not just a summary, but even only the meaning and the more general lines is certainly rather risky. This is therefore a warning to the reader.

In the introductory section, Bhagavan makes us present at the rising of the moral crisis in man: doubt, anguish, the problem of certainty, that of the I and the not-I, and finally the anxious searching for a firm point, of a concept that can take account of the chaotic and contradictory material of experience. Various solutions proposed by the East and the West are little by little undertaken and then criticized and made to pass into other until it reaches an understanding of an opposition as the final foundation: on the one hand, the experience of multiplicity, the changeable, and the illusory refers, as in its condition, to something simple, firm, self-evident: to the I; on the other hand, the I does not have at all a pure experience of itself as such, but is inseparably connected to this dispersed, changeable, and uncertain matter of sensible reality, to which it is opposed. Having reached such a result, it remains to understand its meaning. Here the author moves onto a deductive procedure, discussing the Pranava Vada, i.e., the ``Science of the Sacred Word''.

The central point is that the truly absolute is not that which is simply itself in an emptiness of abstract identity, but rather that which is itself as negation of all that it is not, and that therefore it is at the same time the creative function of its ``other''---of the particular, the contingent, and the changeable---if only as the simple antithesis, against which it strikes the awareness of the universal, if only as the simple matter generated only by giving life to the flame that, devouring it, reaffirms the point of infinity. One therefore has a dialectic concept of the Absolute, enclosed in the sacred word Aum, whose meaning is: aham etat na = ``I-this-not'' = ''I am not this''; i.e., the idea of the Absolute as that supreme simplicity that draws itself from the synthesis of three movements: of an I, of a not-I or ``this'' opposed to it, and of a power that connects the three terms together according to a double function of creation (the I that poses the not-I) and of negation (the not-I dissolved and taken back into the I).

Brahman is exactly this syntheses in its aspect of actual, eternal, simultaneous unity of the three moments; a unity that is neither the one opposed nor the other, neither affirmation nor negation, neither universal nor particular, neither rest not motion, etc. but something that transcends and contains them all, all contradictions and higher reconciliations. It is the \emph{dvandvatitam}, opposed to the \emph{dvandam}, i.e., to the relative, the process of which one can say, like Aristotle, that it is the \emph{act} and the eternal substrate.

As to the \emph{dvandvam}, it therefore includes everything that is determined through relationships, and is expressed in the three moments just mentioned.

\subsection*{Second part}


\paragraph{Pratyagatma}


Pratyagatma, would be the I or Brahman understood in his immanent aspect so that it is determined by its opposite or not-I. It is that which stays fixed and one within every variation, the universal, the point without extension, and the eternal required through the possibility of the experience of everything that is spatially multiple and temporal. From its relation to the not-I, which is essential to it, it nevertheless takes in three determinations:

\begin{enumerate}


\item Jnana = knowledge, as the moment according to which it places for itself and against itself or represents the not-I, the ``this''

\item Kriya = action, as the moment according to which it is affirmed in the not-I

\item Iccha = will or desire according to which the possibility required by the preceding moments, of self-identification of disidentification from the not-I
\end{enumerate}


\paragraph{Mulaprakriti}


Mulaprakriti, matter or not-I. As opposed to the I, which is the universal, it would be pure, dispersed particularity. Nevertheless, in as much as it is also tied to the I, it reflects in itself its characters, i.e., infinity, immobility, and eternity, however as that false infinity, immobility, and eternity that is characteristic to the infinite number of particulars, to the eternal of the ``always finite'' and the ``always changeable''. The way that Bhagavan Das deduces the becoming of nature is interesting. If one could take the not-I in its totality, it would be possible to conceive its negation in a single simultaneous act; but the character of the not-I is instead particularity. Now the particular does not have a \emph{contradictory} that negates it suddenly, but rather a \emph{contrary} that negates it successively. More clearly: being can be negated suddenly by non-being; but heat has instead is negation in the cold that follows it. So the ``I am not this'' of the supreme synthesis, the reflection in mulaprakriti, has the meaning of ``I am not this, but this other; not this other, but still this other'' and so on, in a succession and an indefinite recurrence of opposites that negate each other, the one after the other, in a current through eternal revulsion and eternal craving. Such is the truth of nature, the eternal wheel of beings and rebirths. The unlimitedness of the world as becoming is therefore only a self-reflection, onto the many things, of the unlimitedness of Brahman---the transcendent One without limitation.

\paragraph{Mayashakti}


Mayashakti, which is the demiurgic power from which creation and destruction proceeds and from which the various relationships between I and not-I are determined and governed. Under a first aspect is pravritti, i.e., the direction by which knowledge pushes itself in a certain way beyond itself and, identifying itself, associating itself with a particular material, becomes conscious of a such and such determined finite being. Under a second aspect is nivritti, that through which the I reconverges, detaches itself from particularity, negates it, and little by little elevates itself in every higher degrees of synthetic consciousness and of universality, reflected by the hierarchy of the beings of nature approaching man, to whom in his turn the way of the yogis and rishis is destined up to the limit of the brahmanjanis, or ``those who know'' Brahman.

\paragraph{Conclusion}


Here we cannot even sketch the ample and complex developments given by Bhagavan Das to this conception; on the basis of the ternary rhythm contained in AUM transposed on various ``octaves'', he has the means of giving deductively not only the general lines of a logic, of a music, of an anthropology, of a theory of knowledge, but also of the rituals, symbols, of traditional mythology, thus reaching a true system of Indian knowledge. We pause only on one point. The I in itself, contains and negates in a single act the totality contradicted by the infinite particulars of infinite oppositions. This synthesis therefore is obscured at the level of finite existence, which instead experiences a particular or group of particulars at the same time: these are for them the real and the actual, the rest are pushed back into the past and into the future, as virtual, whose unity (i.e., that total synthesis described above) is nevertheless ideally presupposed by him, in the sense that it makes those universal concepts possible, which is needed in order to be orientated an organize his own finite experience. Now this problem arises: how can one conceive the coexistence of the finite I (jiva) in whom only a part of possible experience is actual, with that ``I'' in which all is actual (Brahman)?

Bhagavan Das' answer is that whatever is virtual for a single particular consciousness, is on the other hand actual for other consciousnesses, and is explained with a rather nice image. Let us suppose a gallery of paintings immersed in the darkness, and a person with a lamp capable of illuminating only one of them at a time; this person will pass from painting to painting, excluding thereby the rest of them, relegating them into a past or a future. Let us now assume as many persons as paintings: for each one of them the same could be said, however the opposition of an actual and a virtual; but in their wholeness they will succeed in entirely illuminating the gallery that they walk though freely substituting themselves ecah one with the other. Such is the relationship of pratyagatma or universal consciousness with the consciousness of the individuals; a pure brainwave pulling itself from the infinite process with which infinite individualities, by identifying and disidentifying themselves, experience, and transcend step by step, they pose and negate the infinite multiplicity characteristic of the non-I.

Even by this fleeting review, the interest that the works here mentioned should generate is obvious. These works should therefore not be missing from the library of all those who turn to the East not out of extrinsic curiosity or a craving for exoticism, but through a true interest; and that therefore they desire to have the spirit and the concept of traditional wisdom, not erudite and philological representations, discourses, and disquisitions about them, as unfortunately is the case of the great number of the publications of our orientalists which, to tell the truth, too often resemble eyeglasses without eyes.

\subsection*{The Idea of the World in the West}
\begin{quotex}
This the the concluding part of review by \textbf{Julius Evola} that originally appeared in Bilychnis, volume XXVI, September, 1925. In this section, he discusses the book \textit{L'idea e il mondo} by R. Pavese, while comparing it to the ideas previously discussed about Bhagavan Das. I know nothing of Pavese and he does not seem to have made any impact on subsequent philosophy. Nevertheless, the review brings up good points. In particular, there is the insistence that speculative thought is vague and useless without a spiritual practice to bring such thought into realization.
\end{quotex}

Regarding what we just said about the Pranavavada\footnote{\url{https://www.makara.us/05ref/01books/pranavavada/pv_toc.htm}}, the reader will have observed by himself---as Bhagavan Das did not fail to observe---the many analogous points that oriental views share with what has been developed in recent European speculative thought. One would almost be tempted to think of a sort of great cycle of culture, at whose end some themes proposed at the beginning would be reintroduced, although on another octave. It is, for example, surprising how close to the Pranavavada is the conception constructed by R. Pavese. In his own way, starting from gnoseological idealism, as opposed to which he reaffirms the legitimate requirement that it is not sufficient to return to the so-called ``transcendental I'' and to a priori logic, but rather it is necessary to determine a conception of these, which takes into account our concrete experience, reality, and physical discoveries in which it has developed. Otherwise, one might just as well exile philosophy into a world of vague and useless abstractions.

The meaning of the conception he proposed is, approximately, the following. God cannot be conceived as absolute transcendence unless in reference to a contrasting, absolute immanence, as eternity is in contrast to temporal becoming. Therefore by showing effectively the process that generates this opposition in order to then reconcile it in a higher inspiration, we have approximately the same idea condensed in the sacred word AUM. Although for Pavese, one pole is the idea as ``logical comprehensiveness'', the other is matter as physical extension and energetic field. The two poles correspond to each other (that which ideally is concept logically comprehending the various known things, objectively is the ``nucleus'' that dominates a ``certain energetic field'', and vice versa) and they convert into each other. It is like a cyclic flux of ``precipitation'' and ``sublimation''. On the one hand there is the idea that realizes itself and, in that way, engenders physical extension, existence in space and time, from that logical comprehensiveness: such is the ``practical'' or natural phase that implies a beginning of relative unconsciousness, since the idea is immersed in it, it loses itself in its act, making itself brute power.

On the other hand, there is the nature that reconverges, the idea that reclaims itself from that multiplicity and that existence into which it had degraded itself, and it reincludes it in a new spiritual unity, that is proper to the concept and knowledge in general. Such is the conscious, ``theoretical'' phase opposed to the natural unconscious process. But the two moments, as was said, are simultaneous and recurrent: it is a flux that rises from one place and descends from the other without interruption, according to an unending vortex that makes spirit from matter and matter from spirit. The two directions, intersecting each other, give rise to nodes or ``inflection points'', in which the eternal cycle is reflected in miniature and which would correspond to individual beings. The identity, the absolute adequation of idea and reality, of comprehension and extension, nucleus and field, simple form and plural matter, is found only in God, the eternal ``concept'' of cosmic becoming, to become what is His ``conception'' or content, in which He circulates with alternating rhythms. In finite individuals, identity is instead a ``having to be'', something to which they tend, since the order of nature is precisely defined by an eternal imbalance of the two moments from which its becoming arises.

Pavese's originality lies in his bold placing in reciprocal conversions the logical comprehension of a concept with the extension of an energetic field understood as its objective correlative. That gives him a way to graft onto a dialectic scheme, that reminds us of Proclus and Spinoza, the views of recent physics and mechanics, and thereby extends the speculative penetration right into the very heart of the concreteness established in science. In the second place, in the effort of dynamizing, of making ``fluid'', of gathering in simultaneity the moments of the process, it results in a type of classificatory system of other conceptions. Opposed to that, there is the lack of an organic and systematic exposition and a certain obscurity and convolutedness of expression.

In similar oriental and occidental conceptions, one must nevertheless ask up to what point is it critically justifiable to place oneself at the Absolute, in order to deduce all the rest from that. A Cartesian argument attempted by Pavese is, to tell the truth, inadequate. Similarly for Bhagavan Das, it is true that he first proceeds inductively starting from the analysis of the nature of human experience; but when he then tries to explain the result of that analysis with a deduction, even he makes, and cannot help but make, a leap: the two methods cannot be joined together. Deduction can offer only probable explicative principles in order to interpret the data of fact and, as such, they are deprived of necessity and convertibility. It is therefore necessary to note that the criterion of certainty for the Oriental is not criticism and logic, but rather that of a ``transcendental experience''. It is, for example, said that the Vedas are the expositions of what the \emph{rishis} have \emph{seen}. For the \emph{rishi}, meaning those who have realized themselves up to the level of Brahman, the deductive procedure of their knowledge could have in itself the character of intrinsic evidence and a justification, that cannot appear to those who look at it from the level of finite existence.




\flrightit{Posted on 2013-04-28 by Cologero}

